\documentclass[10pt]{article}
\usepackage{../../../estilo/cartoes}

\hypersetup{
  pdftitle={Vanguarda 44: Frente Leste -- Cartao de Cenario},
  pdfauthor={Felipe Ramires Terrazas},
}

\begin{document}

% ------------------------------------------------------------------
% Folha 1 -- Cartao 1
% ------------------------------------------------------------------
\begin{cardpage}

\cardcell{0}{297}{1}{STALINGRADO}{
\textbf{Ago.\ 1942--fev.\ 1943, Stalingrado, União Soviética.}

\textbf{Forças:} Alemães (Atacante) vs.\ Soviéticos (Defensor).

\textbf{Qualidade de Tropa:} Alemães Veteranos; Soviéticos Regulares (compensam com o bônus de CQC da doutrina nacional).

\textbf{Regra Especial ``Não Há Terra Além do Volga'':} unidades soviéticas nunca recuam voluntariamente: a ordem Correr não pode ser usada para se afastar do lado alemão da mesa (só para flanquear ou avançar).

\textbf{Objetivo:} Soviéticos vencem se ainda controlarem ao menos 1 de 2 a 3 edifícios-chave definidos na mesa até o turno 10. Alemães vencem se tomarem todos.

\textbf{Duração:} 10 turnos.
}

\end{cardpage}

\end{document}
